\section{Schematicity of abelian torsors and their Picard functors}
\label{sec:abelian-torsor-schemes}

An abelian scheme in this section is a smooth proper group scheme of
finite presentation with geometrically connected fibres. Projectivity is
not included in the definition. All torsors and contracted products are
for the fppf topology.

The distinction between an algebraic space and a scheme matters for
torsors under abelian schemes. We prove that changing the structure group
preserves schematicity when the original torsor is a scheme.

This section proves Theorem~\ref{thm:schematic-abelian-pushout} and Corollary~\ref{cor:picard-schematic-abelian-torsor}.

There is no assumption on the kernel of $h$. In particular, the theorem
applies when the kernel has positive dimension. The proof combines
Raynaud's criterion for effective descent with the affine-neighborhood
and ampleness results for connected group actions in
\cite[\S6.6, Theorem~2(c),(d), pp.~164--166]{BLR}.
We give the finite-cover and approximation steps explicitly.

Call a scheme \emph{AF} if every finite set of its points is contained
in an affine open subscheme. We use the following form of Raynaud's
criterion \cite[Theorem~4.2, pp.~5-16--5-18]{RaynaudDescent}.
Let $F$ be an fpqc sheaf over $S$. Suppose that an fpqc cover $V\to S$
makes $F_V$ a locally noetherian scheme, and that a finite surjective
map $W\to S$ makes $F_W$ a scheme. Set
\[
 Y=F_W,\qquad Z=F_{W\times_S W},
\]
and let $v_1,v_2:Z\rightrightarrows Y$ be the projections.
Then $F$ is a scheme if and only if
$v_2(v_1^{-1}(y))$ is contained in an affine open of $Y$ for every
$y\in Y$. Thus an AF scheme $Y$ suffices. The finite cover $W\to S$
is not required to be flat, and the criterion does not require a
descended line bundle.

\begin{lemma}\label{lem:semilocal-multisection}
Let $R$ be a noetherian semilocal ring. Let $Y\to\Spec R$ be projective,
smooth and surjective, of pure relative dimension $d$.
There is a closed subscheme $Z\subset Y$ that is finite locally free
and surjective over $\Spec R$.
\end{lemma}

\begin{proof}
For $d=0$ take $Z=Y$. Suppose $d>0$, and choose a very ample line
bundle $L$ on $Y$. Write $\mathfrak m_1,\ldots,\mathfrak m_r$ for the
maximal ideals of $R$ and $Y_i$ for the corresponding fibres.
Graded Noether normalization on each $Y_i$ supplies $d$ homogeneous
sections of positive powers of $L_i$ whose common zero scheme is
nonempty and zero-dimensional: take $d$ of the $d+1$ homogeneous
parameters of a finite map to $\PP^d_{\kappa(\mathfrak m_i)}$.
This construction works over finite residue fields as well.
Taking powers puts all the sections in a common degree without
changing their common zero sets.

Set $I=\bigcap_i\mathfrak m_i$. For sufficiently large $n$, Serre
vanishing gives
\[
 H^1(Y,I\OO_Y\otimes L^n)=0.
\]
After taking further common powers, the chosen sections on the disjoint
union of the $Y_i$ therefore lift to $d$ sections of $L^n$ on $Y$.
Let $Z$ be their common zero scheme. It is projective, and its closed
fibres are nonempty and zero-dimensional.

The locus where a fibre of $Z\to\Spec R$ has positive local dimension
is closed in $Z$ by \cite[IV$_3$, 13.1.3]{EGA}. Its image is closed
by properness and contains no closed point of $\Spec R$, so is empty.
Thus $Z\to\Spec R$ is proper and quasi-finite, hence finite.

At each point above a maximal ideal of $R$, the chosen fibre sections
form a system of parameters in a regular local ring of dimension $d$.
They form a regular sequence. Since $Y$ is $R$-flat, the fibrewise
regular-sequence criterion \cite[IV$_3$, 11.3.8]{EGA} makes $Z$ flat
at those points. The flat locus is open in $Z$. The image of its
closed complement under the finite map to $\Spec R$ avoids every
closed point, so that complement is empty. Therefore $Z$ is finite
locally free over $R$. Its closed fibres are nonempty, and its rank
is locally constant; the map is consequently surjective.
\end{proof}

\begin{proof}[Proof of Theorem~\ref{thm:schematic-abelian-pushout}]
Write $Q=h_*T$. It is a $B$-torsor algebraic space, and
$Q_T\simeq B_T$. We first work over an excellent noetherian local
scheme $S=\Spec R$.

Choose an affine open $U\subset T$ meeting the closed fibre.
The map $T\to S$ is smooth, so the image of $U$ is open and contains
the closed point of $S$; it is all of $S$. The $A$-action is transitive
on geometric fibres, whence $A\cdot U=T$.
The affine-neighborhood theorem \cite[\S6.6, Theorem~2(c)]{BLR}
makes $T$ an AF scheme. The same argument, starting with an affine
neighborhood of the identity in $B$, makes $B$ AF.

Let $W\to S$ be the normalization of $S_{\mathrm{red}}$.
It is finite and surjective because $S$ is excellent; $W$ is normal,
affine and semilocal. A finite morphism preserves AF on its source:
put the finitely many image points in an affine open and take its
affine inverse image. Thus $T_W$ and $B_W$ are AF.

Choose one point of $T_W$ above each closed point of $W$.
A single affine open contains all these points. Its image in $W$
is open by smoothness, contains every closed point and hence equals
$W$. Its saturation under $A_W$ is all of $T_W$.
Normality of $W$ and \cite[\S6.6, Theorem~2(d)]{BLR} now supply
a relatively ample line bundle on $T_W$; therefore $T_W$ is
projective over $W$. The same argument makes $B_W$ projective.
The use of one affine open meeting all closed fibres is essential.

By Lemma~\ref{lem:semilocal-multisection}, there is a finite locally
free surjective closed subscheme $Z\subset T_W$. Its inclusion gives
a section of $T_Z$, and thus $Q_Z\simeq B_Z$.
Since $B_W$ is projective and $Z$ is affine, $B_Z$ is AF.
Apply Raynaud's criterion over $W$, using $Z\to W$ both as the fpqc
cover and as the finite cover. It follows first that $Q_W$ is a
scheme.

We next prove projectivity of this scheme. The map
$q:Q_Z\to Q_W$ is finite locally free and surjective, and $Q_Z$ is
projective over the affine scheme $W$. Choose an ample line bundle
$L$ on $Q_Z$. Its norm
\[
 \operatorname{Nm}_{q}(L)
 =\det(q_*L)\otimes\det(q_*\OO_{Q_Z})^{-1}
\]
is an invertible sheaf on $Q_W$, and is ample by
\cite[II, Proposition~6.6.1 and Corollary~6.6.4]{EGA}.
Properness of $Q_W\to W$ follows by fpqc descent from the $B$-torsor
structure. Hence $Q_W$ is projective over $W$, and in particular AF.
The norm is used only after the target has been proved to be a scheme.

Apply Raynaud's criterion a second time, over $S$.
The fpqc cover is $T\to S$, on which $Q_T\simeq B_T$ is a locally
noetherian scheme. The separate finite surjective cover is $W\to S$,
on which $Q_W$ is AF. The criterion proves that $Q$ is a scheme.
No flatness of $W\to S$ has been used.

We justify passage from local rings to open neighborhoods.
Start with an excellent noetherian affine $S$ and $s\in S$, and
form its finite normalization $W\to S$ before localizing.
Over $W_s=W\times_S\Spec\OO_{S,s}$ the preceding argument constructs
projective embeddings of $T_{W_s}$ and $B_{W_s}$, and a
finite locally free surjective closed subscheme $Z_s\subset T_{W_s}$.
As $W_s$ is the inverse limit of the $W_U$ for affine neighborhoods
$U$ of $s$, these finitely presented objects and maps descend to
some $W_U$. After shrinking $U$, the embeddings remain closed
immersions and $Z_U\to W_U$ is finite locally free and surjective.
These are finite-presentation limit properties
\cite[IV$_3$, 8.8--8.10 and 11.2.6(ii)]{EGA}.
The two applications of Raynaud's criterion then take place over
the open neighborhood $U$. Thus $Q_U$ is a scheme. Such
neighborhoods cover $S$.

Finally let $S$ be arbitrary, and work on an affine part of it.
The schemes $A,B,T$, their group laws, the action, the homomorphism
$h$, and the inverse of the torsor isomorphism
$A\times_S T\simeq T\times_S T$ are finite-presentation data.
They descend to a model over a finitely generated $\Z$-algebra.
After increasing the model, the finitely many identities hold,
and properness, smoothness and surjectivity hold as well
\cite[IV$_3$, 8.10.5; IV$_4$, 17.7.8(ii)]{EGA}.
For a proper smooth morphism the number of geometric connected
components is locally constant \cite[IV$_3$, 15.5.9(ii)]{EGA};
restrict the model to the open and closed locus through which
the original base factors. The fibres are then geometrically
connected. The resulting base is excellent and noetherian.
The pushout is compatible with base change, so the result there
proves schematicity over the original base.
\end{proof}



The conclusion over a connected locally noetherian base is stated
in \cite[\S1.7.2, p.~13]{ACMPrym}. The proof below spells out the varying Néron--Severi
parameter and the possibly nontrivial torsor of line-bundle classes
on $A$, then applies Theorem~\ref{thm:schematic-abelian-pushout}.
The contracted-product identity is classical; see
\cite[Theorem~3.0.3, p.~664]{Alexeev}.

\begin{proof}[Proof of Corollary~\ref{cor:picard-schematic-abelian-torsor}]
The same approximation as in the theorem reduces to a noetherian
base, since relative Picard sheaves commute with arbitrary base
change. Write $P_A=\Pic_{A/S}$ and $B=A^\vee$.
The Picard functor is a separated algebraic space locally of finite
presentation \cite[Section~108.9, Lemmas~3--4]{Stacks}:
the proper smooth morphism has geometrically integral fibres and
its degree-zero direct image is universally the structure sheaf.
Its open neutral subfunctor $B$ is an abelian scheme; see,
with the definition of abelian scheme used here,
\cite[Corollary~4.6, author version pp.~14--15]{LaurentSchroer}.
Since $B$ is proper over $S$ and $P_A$ is separated, $B$ is also
closed in $P_A$.

Form the quotient algebraic space $N=P_A/B$ for the free action
of the smooth group $B$ \cite[Tag~06PH]{Stacks}. The map $P_A\to N$ is a $B_N$-torsor.
Local finite presentation of $N\to S$ descends through this smooth
surjection. The pullback of $\Delta_{N/S}$ to $P_A\times_S P_A$
is the inverse image of $B$ under the difference map, so is closed.
Thus $N$ is separated over $S$. Its geometric fibres are discrete,
being the quotients of the geometric Picard groups by their open
identity components. Hence $N\to S$ is locally quasi-finite, and
the separated locally quasi-finite algebraic-space theorem
\cite[Tag~03XX]{Stacks} makes $N$ a scheme. We do not assert
that $N$ is flat or that its ranks are locally constant.

Regard $Q_A=P_A\to N$ as the $B_N$-torsor of the universal class
$\delta\in N(N)$. Inversion $\iota=[-1]^*$ acts by $-1$ on $B$
and by $+1$ on $N$
\cite[\S\S3.1--3.2, pp.~1311--1312]{PR}.
Its fixed subsheaf $R\subset Q_A$ is a $B[2]_N$-torsor:
if $x$ is a local point of $Q_A$ and $d=\iota(x)-x\in B$, then
\[
 \iota(x+b)-(x+b)=d-2b.
\]
The finite locally free surjection $[2]:B\to B$ therefore gives
fixed points locally, whose differences lie in $B[2]$.
Consequently $R\to N$ is finite locally free and surjective, and
$R$ is a scheme. It supplies a tautological section of $(Q_A)_R$,
so $(Q_A)_R\simeq B_R$.

For each affine open $V\subset N$, the scheme $R_V$ is affine and
noetherian. An abelian scheme over an affine base is AF
\cite[application after Corollary~3.8, p.~5-15]{RaynaudDescent}.
Raynaud's criterion, using $R_V\to V$ for both covers, makes
$(Q_A)_V$ a scheme. Hence $Q_A$ is a scheme over $N$.
This proves schematicity of the component torsor; it does not
replace that torsor by the trivial one.

For a local section $t$ of $T$, let
$j_t:A\to T$ send $a$ to $t+a$.
The induced identifications of $\Pic_{T/S}$ with $P_A$ identify
the neutral parts with $B$. Changing $t$ acts by translation on
$A$, and translation fixes both $B$ and $N$. The quotient parameters
therefore glue to a canonical map
\[
 P_T=\Pic_{T/S}\longrightarrow N,
\]
which is a $B_N$-torsor. The universal class $\delta$ defines
a homomorphism
\[
 \lambda_\delta:A_N\longrightarrow B_N,\qquad
 a\longmapsto \tau_a^*L\otimes L^{-1},
\]
where $L$ is any local representative of $\delta$.
The expression is independent of that representative because
translation fixes $B$, and is a homomorphism by the theorem of
the square.

The scheme $U=T_N\times_N Q_A$ is a torsor under the abelian
$N$-scheme $A_N\times_N B_N$. On relative Picard classes define
\[
 m:U\longrightarrow P_T,\qquad
 (t,L)\longmapsto(j_t^{-1})^*L.
\]
There is no need to choose a universal line bundle in this formula.
Since $j_{t+a}^{-1}=\tau_{-a}\circ j_t^{-1}$, one has
\[
 m(t+a,L+b)=m(t,L)+b-\lambda_\delta(a).
\]
Thus $m$ is equivariant for
$h(a,b)=b-\lambda_\delta(a)$ and induces an isomorphism
$h_*U\simeq P_T$ of $B_N$-torsors. It is an isomorphism because,
after choosing $t$ and $L$ fppf locally, it is the identity of $B_N$.
Theorem~\ref{thm:schematic-abelian-pushout} now makes $P_T$ a
scheme. Its separatedness and local finite presentation were
already properties of the Picard algebraic space.
\end{proof}

\begin{remark}
These results settle the schematic abelian-torsor subcase of the
question whether a proper flat family with geometrically normal
or geometrically integral fibres has a schematic full Picard functor.
They do not settle that question for general higher-dimensional
families. In particular, geometric normality alone does not imply
that the whole neutral Picard space is smooth or flat.
\end{remark}
